\documentclass[11pt,a4paper]{article}
\usepackage[utf8]{inputenc}
\usepackage[T1]{fontenc}
\usepackage{graphicx}
\usepackage{booktabs}
\usepackage{hyperref}
\usepackage{amsmath}
\usepackage[margin=1in]{geometry}
\usepackage{natbib}
\usepackage{textcomp}

\title{Presentation Report: Dispersion, Not Displacement --- Loss of Alveolar Epithelial State Coherence in Lethal COVID-19}
\author{Chimdi Walter Ndubuisi and Pierce Taylor\\[0.5em]
\textit{Instructor: Dr.\ Toni Kazic}}
\date{May 12, 2026}

\begin{document}
\maketitle

\begin{abstract}
The alveolar epithelium repairs itself through AT2-to-AT1 differentiation, but
this process is overwhelmed in lethal COVID-19. Two geometric models could
describe the resulting injury: coherent displacement of cells toward a single
damaged state, or heterogeneous dispersion across multiple injury states. We
tested these models using 22,128 alveolar epithelial cells from the Columbia/NYP
COVID-19 Lung Atlas (SCP1219; 7 healthy donors, 20 COVID-19 donors), scoring
eight gene programs along diffusion pseudotime and evaluating robustness through
ten pre-specified ablation analyses, Harmony batch correction, an orthogonal
trajectory algorithm (Palantir), and independent replication on 89,736 cells from
618 donors across 35 datasets.

Coherent centroid displacement was not supported (one-sided Mann--Whitney
$p = 1.0$; rank-biserial $r = 0.14$ opposite to the hypothesized direction).
Within-group dispersion was strongly supported (variance ratio 2.3; Levene
$p < 10^{-30}$) and replicated decisively cross-cohort (variance ratio 1.65;
$p \approx 10^{-137}$). COVID-19 cells were enriched at higher pseudotime
(permutation $p = 0.001$; median difference $+0.048$ on a $[0,1]$ axis), a
signal preserved in 27/27 leave-one-donor-out iterations, amplified 2.6-fold
under Harmony batch correction, confirmed at reduced magnitude by Palantir
($+0.028$), and replicated at the donor-pseudobulk level in the independent
cohort ($p = 0.002$). A KRT8$^+$/CLDN4$^+$ transitional population was enriched
3.1-fold in the primary atlas but did not transfer cross-cohort under
marker-threshold definitions. The fine-grained ordering of injury programs along
pseudotime was sensitive to gene-list composition and is treated as
hypothesis-generating.

These results support a model in which lethal COVID-19 produces
dispersion-dominated loss of alveolar epithelial state coherence rather than a
coherent directional shift. Within-group dispersion and donor-level trajectory
direction are the most portable findings across cohorts and methods;
annotation-dependent and threshold-based measures are atlas-specific.
\end{abstract}

\tableofcontents
\newpage

\section{Introduction}

The alveolar epithelium is the functional unit of gas exchange in the lung. Two
cell types sustain this function: AT1 cells, which form the thin gas-exchange
barrier across approximately 95\% of the alveolar surface, and AT2 cells, which
secrete pulmonary surfactant and serve as the resident stem cells of the alveolar
compartment. When AT1 cells are lost to injury, AT2 cells proliferate and
differentiate through a KRT8$^+$/CLDN4$^+$/SFN$^+$ transitional state to
regenerate the gas-exchange surface \citep{Barkauskas2013, Kobayashi2020, Strunz2020}.
This AT2-to-AT1 differentiation axis is the primary mechanism of alveolar
self-repair.

In lethal COVID-19, this regenerative system is overwhelmed. SARS-CoV-2 enters
AT2 cells via ACE2, directly compromising both the surfactant system and the
regenerative progenitor pool \citep{Ziegler2020, Hou2020}. A dysregulated
inflammatory response inflicts further collateral damage. Autopsy studies describe
diffuse alveolar damage, loss of normal architecture, and accumulation of
transitional epithelial cells that are neither fully AT2 nor fully AT1 \citep{Borczuk2020, Melms2021}. Single-cell atlases of these tissues have revealed cell-type
compositional changes and differentially expressed genes \citep{Melms2021, Delorey2021}, but the cluster-based approach they employ forces continuous biology into
discrete bins.

A trajectory-based approach offers a more informative alternative. By treating
each cell's expression profile as a position on a continuous transcriptomic
manifold and ordering cells along pseudotime \citep{Haghverdi2016, Setty2019}, one
can ask whether the molecular evidence is consistent with staged failure and
characterize the geometry of the injury pattern. However, this approach raises a
question that is rarely addressed explicitly: what is the expected \textit{shape} of
epithelial injury in transcriptomic space?

Two competing geometric models make distinct predictions (Figure~\ref{fig:manifold}). Under a
\textbf{coherent displacement model}, injury shifts the entire cell population away from
the healthy centroid toward a single new location in gene-expression space,
implying a uniform damaged state. Under a \textbf{dispersion model}, injury fans cells
outward from the healthy region into a heterogeneous cloud of partially
overlapping injury states, increasing within-group variance without necessarily
moving the centroid. These models make opposite predictions for two standard
metrics: a centroid-displacement test (significant under displacement, null under
dispersion) and a variance test (null under displacement, significant under
dispersion). Both models predict enrichment at higher pseudotime values under a
root-in-healthy-cells design.

Here, we test these models using alveolar epithelial cells from the Columbia/NYP
COVID-19 Lung Atlas \citep{Melms2021}. We pre-specified three statistical criteria ---
centroid displacement, within-group dispersion, and pseudotime enrichment --- and
evaluated robustness through ten ablation analyses, Harmony batch correction, an
orthogonal trajectory algorithm (Palantir), and independent cross-cohort
replication on 89,736 alveolar cells from 618 donors across 35 datasets. The
result is an explicit accounting of which conclusions are robust across
methodological choices, which are portable across cohorts, and which are
atlas-specific.

\section{Results}

\subsection{Cohort composition and alveolar subset validation}

We extracted 22,128 alveolar epithelial cells from SCP1219 by retaining cells
annotated as AT1, AT2, or ECM-high epithelial in the original atlas \citep{Melms2021}:
9,608 AT1 cells, 11,341 AT2 cells, and 1,179 ECM-high transitional cells from 27
donors (12,181 cells from 7 healthy controls; 9,947 from 20 COVID-19 cases).
Per-donor alveolar cell counts ranged from 62 to 2,850, and 10 of 27 donors
contributed fewer than 300 cells (Supplementary Table S5). The true inferential
sample size is therefore 27 donors (7 versus 20), and we rely on donor-level
analyses and leave-one-donor-out ablation throughout to ensure that no single
donor drives any primary finding.

Marker-gene expression confirmed annotation quality: AT2 cells expressed \textit{SFTPC}
and \textit{SFTPA1}; AT1 cells expressed \textit{AGER} and \textit{HOPX}; ECM-high transitional cells
expressed \textit{KRT8} and \textit{CLDN4} at elevated levels relative to both AT1 and AT2
populations (Supplementary Table S3). The marker profile of the ECM-high
compartment is consistent with the KRT8$^+$/CLDN4$^+$ damage-associated transient
progenitor (DATP) population described in mouse bleomycin injury and human
idiopathic pulmonary fibrosis \citep{Kobayashi2020, Strunz2020, Adams2020, Habermann2020}.

\begin{figure}[htbp]
\centering
\includegraphics[width=\textwidth]{results/figures/fig2_manifold.png}
\caption{\textbf{Alveolar epithelial manifold.} UMAP of 22,128 alveolar epithelial cells colored by condition (left), cell type annotation (center), and diffusion pseudotime (right). Pseudotime is rooted at the healthy AT2 centroid and normalized to [0, 1].}
\label{fig:manifold}
\end{figure}

\subsection{Two competing models of epithelial injury: centroid displacement fails, dispersion succeeds}

The central question of this study is whether alveolar epithelial injury in
lethal COVID-19 takes the form of coherent displacement or heterogeneous
dispersion. We tested both using the same cells in the same 50-dimensional PCA
space.

\subsubsection{Centroid displacement was not supported}

We computed the Euclidean distance from each cell to the centroid of healthy cells
in PCA space and tested whether COVID-19 cells were farther from this centroid
than healthy cells (one-sided Mann--Whitney U, alternative COVID $>$ Healthy). The
test was not significant ($U = 5.206 \times 10^{7}$, $p = 1.0$; rank-biserial
$r = 0.14$ pointing opposite to the alternative). The median distance was 10.68
for COVID-19 cells and 11.56 for healthy cells --- COVID-19 cells were, if
anything, \textit{closer} to the healthy centroid than healthy cells were to their own
centroid. This result was pre-specified as one of three core criteria and the data
do not support it.

\subsubsection{Within-group dispersion was strongly supported}

For each condition we computed each cell's Euclidean distance to that condition's
own centroid in PCA space. The variance of these distances was 33.28 in COVID-19
and 14.58 in controls (ratio 2.3). Levene's test yielded a statistic of 276.5
with $p < 10^{-30}$ (Table~\ref{tab:three-tests}). COVID-19 alveolar cells occupy a markedly broader
region of transcriptomic space than healthy cells, consistent with fragmentation
into heterogeneous injury states rather than coherent movement to a single new
state.

\subsubsection{Interpretation: dispersion, not displacement}

These two results together discriminate between the competing models. The absence
of centroid displacement rules out the simpler model in which all COVID-19 cells
shift coherently toward one damaged phenotype. The strong dispersion signal
supports the alternative model in which injury fans cells outward across multiple
partially overlapping states while some cells remain near homeostatic identity. A
test that measures mean displacement averages over the near-healthy and the
distant subpopulations and detects neither. A test that measures within-group
variance captures exactly this structure. We propose that future studies of
epithelial state loss should adopt dispersion-based or trajectory-based metrics in
preference to centroid-displacement tests, which implicitly assume a unimodal
shift that injury biology does not produce.

\subsection{Dispersion-dominated loss of state coherence is the most portable finding}

The dispersion result is the strongest and most portable finding in this study. In
the primary atlas, the 2.3-fold variance ratio between COVID-19 and healthy cells
is the largest effect size among all tested metrics.

To evaluate cross-cohort portability, we constructed an independent replication
cohort of 89,736 alveolar AT1/AT2 nuclei from the cellxgene census (35 datasets,
618 donors; the Melms/SCP1219 dataset was excluded by its identifier). We scored
the same eight gene programs and defined an injury composite as the mean of the
six non-homeostatic program scores. Levene's test on the injury composite yielded
a statistic of 626.2 with $p \approx 10^{-137}$ and a variance ratio of 1.65
(COVID variance 0.0505; normal variance 0.0306). The direction and qualitative
magnitude of the primary-cohort dispersion finding replicate decisively.

\subsubsection{Donor-level dispersion}

At the donor level, the pattern is more nuanced. Per-donor within-donor median
distance to own centroid was 7.77 for COVID-19 donors and 10.58 for healthy
donors (Mann--Whitney $p = 0.999$; Levene $p = 0.071$). Individual COVID-19
donors are not more internally dispersed than individual healthy donors; rather,
the cell-level dispersion signal arises from the diversity of states \textit{across}
COVID-19 donors. This is consistent with a model in which each donor's cells
occupy a different region of injury space, collectively producing the fan-out
visible at the cell level. The within-donor dispersion bootstrap 95\% CI for the
difference in medians was [$-30.7$, 16.5], confirming that the cell-level Levene's
test captures between-donor heterogeneity rather than within-donor inflation.

\begin{figure}[htbp]
\centering
\includegraphics[width=\textwidth]{results/figures/fig3_dispersion_centerpiece.png}
\caption{\textbf{Dispersion-dominated loss of state coherence.} (a) Within-group distance-to-centroid distributions for COVID-19 and healthy cells (Levene $p < 10^{-30}$; variance ratio 2.3). (b) Distance-to-healthy-centroid distributions (displacement test; Mann--Whitney $p = 1.0$). (c) Donor-level within-donor dispersion. (d) Replication cohort dispersion (Levene $p \approx 10^{-137}$; variance ratio 1.65). (e) Summary statistics table.}
\label{fig:dispersion}
\end{figure}

\subsection{COVID-19 cells are enriched at higher pseudotime along a donor-robust injury axis}

We constructed a transcriptomic manifold (3,000 HVGs, 50-component PCA,
$k$-nearest-neighbor graph with $k = 15$, 15-component diffusion map) and
computed diffusion pseudotime rooted at the cell nearest the healthy AT2 centroid
in diffusion-component space, min-max normalized to [0, 1].

\subsubsection{Pseudotime shift}

COVID-19 cells were shifted toward higher pseudotime (median 0.111 versus 0.063;
observed difference +0.048). A 1,000-iteration permutation test (condition labels
shuffled, median difference recomputed) placed the null distribution tightly at
zero (mean $6.0 \times 10^{-5}$, std $1.24 \times 10^{-3}$); the observed value
exceeded all permutations ($p = 0.000999$). The shift is modest in absolute
magnitude (approximately 5\% of the normalized range) but statistically robust.

\subsubsection{Donor-level robustness}

Leave-one-donor-out analysis (27 iterations) preserved the direction of the shift
in every iteration, with magnitude ranging from 0.028 to 0.202 and the bulk of
iterations clustering between 0.047 and 0.086. No single donor was necessary for
the pseudotime shift to remain positive.

All 20 COVID-19 donors had median pseudotime $\geq 0.087$, while 5 of 7 control
donors had median pseudotime $< 0.090$. No COVID-19 donor had a median pseudotime
below the lowest healthy-donor median.

\subsubsection{Harmony batch correction amplifies the shift}

Harmony \citep{Korsunsky2019} with \texttt{donor\_id} as batch key converged in three
iterations and yielded a pseudotime shift of +0.125, approximately 2.6-fold
larger than the uncorrected value. This amplification is biologically
interpretable: donor-level batch effects, confounded with condition in this
autopsy-versus-organ-donor cohort, act as within-donor clustering forces in PCA
space. Removing that within-donor clustering while preserving between-condition
structure unmasks a larger condition-associated shift. The fact that batch
correction \textit{strengthens} rather than weakens the signal is strong evidence that
the pseudotime shift is not a donor-effect artifact.

\subsubsection{Palantir cross-validation}

Palantir \citep{Setty2019} on the same manifold and root cell returned a median
pseudotime shift of +0.028 between conditions, of the same sign as diffusion
pseudotime and comparable in magnitude. The program-ordering correlation with the
primary analysis was $\rho = 0.48$. This cross-algorithm agreement on the
direction and approximate magnitude of the shift indicates that the signal is a
property of the underlying manifold rather than an artifact of one algorithm.

\subsubsection{Donor-level replication}

In the independent replication cohort, the cell-level one-sided Mann--Whitney test
on the injury composite returned $p = 1.0$ (consistent with the cell-level
centroid-displacement failure in the primary cohort). However, pseudobulking to
donor-level medians and testing 43 COVID donors versus 575 control donors yielded
$p = 0.002$ in the expected direction. This pattern --- cell-level unimodal-shift
tests fail, donor-level aggregated tests succeed --- mirrors the primary cohort
and reinforces the view that cell-level displacement tests are not the appropriate
instrument for this class of biology.

\begin{figure}[htbp]
\centering
\includegraphics[width=\textwidth]{results/figures/fig4_donor_pseudotime.png}
\caption{\textbf{Donor-level pseudotime analysis.} (a) Pseudotime density distributions by condition. (b) Per-donor median pseudotime (bars), sorted by value. (c) Bootstrap distribution (10,000 replicates) for the difference in donor-level median pseudotime between conditions; 95\% CI [$-0.013$, 0.086]. (d) OLS model summary: coefficient +0.043 ($p = 0.013$; $R^{2} = 0.22$).}
\label{fig:donor-pseudotime}
\end{figure}

\begin{figure}[htbp]
\centering
\includegraphics[width=0.8\textwidth]{results/figures/fig4_pseudotime_density.png}
\caption{\textbf{Pseudotime density by condition.} Distribution of diffusion pseudotime values for COVID-19 (red) and healthy control (blue) alveolar epithelial cells.}
\label{fig:pseudotime-density}
\end{figure}

\subsection{A KRT8$^+$/CLDN4$^+$ transitional compartment is enriched in the primary atlas but is not a portable cross-cohort signature}

The ECM-high transitional population was enriched 3.1-fold in COVID-19 relative
to controls (848/9,947 COVID cells [8.5\%] versus 331/12,181 control cells [2.7\%];
$\chi^{2}$ $p < 10^{-30}$). These cells occupied intermediate pseudotime
positions (mean $\approx 0.3$--$0.5$), expressed \textit{KRT8}, \textit{CLDN4}, and partial AT1
markers (Supplementary Table S3), and are consistent with damage-associated
transient progenitors described in mouse lung injury \citep{Kobayashi2020, Strunz2020}
and aberrant basaloid cells described in human IPF \citep{Adams2020, Habermann2020}.
Their enrichment in COVID-19 and their intermediate pseudotime position are
together consistent with accumulation of cells that initiated AT2-to-AT1 repair
but did not complete it. This interpretation is correlational; we have a
cross-sectional snapshot, not a lineage trace.

\subsubsection{Cross-cohort portability of marker-threshold DATP calling fails}

In the independent replication cohort, we defined DATP-positive cells as those
with simultaneously positive expression of KRT8 and CLDN4 (log-normalized count
$> 0$ for both). The fraction was 22.9\% in COVID-19 versus 47.6\% in controls (fold
enrichment 0.48x; $\chi^{2}$ $p \approx 10^{-228}$), the opposite direction from
the primary atlas. This failure reflects cross-cohort non-equivalence of simple
marker-threshold calls: study-specific normalization, capture efficiency, and
ontology labeling all affect how many cells cross a fixed expression threshold. We
therefore report DATP enrichment as a primary-atlas observation, not a
cross-cohort-validated finding.

\begin{figure}[htbp]
\centering
\includegraphics[width=\textwidth]{results/figures/fig5_transitional.png}
\caption{\textbf{Transitional cell compartment.} (a) UMAP highlighting ECM-high transitional cells. (b) Cell-type composition by condition (stacked bar); 3.1-fold enrichment of transitional cells in COVID-19. (c) Marker gene expression across cell types. (d) Pseudotime density distributions by cell type.}
\label{fig:transitional}
\end{figure}

\subsection{Ablation analyses distinguish robust from fragile conclusions}

Ten pre-specified ablations systematically varied cell-type scope, embedding
method, batch correction, gene-program composition, pseudotime root, trajectory
algorithm, donor identity, and null-model assumptions (Supplementary Table S6).
Each ablation reported three metrics: median pseudotime difference (MPD),
program-ordering Spearman $\rho$ versus the primary analysis, and displacement
rank-biserial $r$.

\subsubsection{Robust findings}

The direction of the pseudotime shift was preserved across all perturbations that
retained the alveolar compartment and a healthy-cell root: 27/27
leave-one-donor-out iterations (MPD range 0.028--0.202); all embedding variants
(MPD = +0.048); the batch-correction fallback (MPD = +0.048); and the
apoptosis-gene removal (MPD = +0.074). The curated-only gene-program variant
produced identical program ordering ($\rho = 1.0$). The displacement effect size
was never significant in the pre-specified direction ($r \leq 0.14$ across all
ablations), confirming that the centroid-displacement failure is a robust feature
of the data, not an artifact of a single pipeline choice.

\subsubsection{Fragile findings}

The fine-grained ordering of injury programs within the upper half of pseudotime
was sensitive to gene-list choice: substituting MSigDB hallmark sets for curated
lists inverted the ordering ($\rho = -0.80$). The pseudotime shift collapsed when
airway epithelial cells were included (Ablation 6, MPD = $-0.008$), confirming that
the trajectory is a property of the alveolar compartment specifically. Rooting
pseudotime in a COVID-extreme cell inverted the program ordering ($\rho = -0.10$),
as expected mathematically.

\subsubsection{Summary}

The existence and direction of the pseudotime shift, and the dispersion finding,
are robust across the full ablation grid. The fine-grained injury-program sequence
and the specificity of the trajectory to alveolar-only analyses are
methodologically sensitive. We report the fine ordering as hypothesis-generating,
not as a confirmed sequence.

\begin{figure}[htbp]
\centering
\includegraphics[width=\textwidth]{results/figures/fig6_ablation_matrix.png}
\caption{\textbf{Ablation robustness matrix.} Ten pre-specified ablations (rows) evaluated across three metrics (columns): median pseudotime difference, program-ordering Spearman $\rho$ versus primary, and displacement rank-biserial $r$. Pseudotime direction is robust across all alveolar-compartment ablations; fine program ordering is sensitive to gene-list choice.}
\label{fig:ablation}
\end{figure}

\subsection{Gene-program dynamics: identity loss is stable, fine ordering is hypothesis-generating}

We scored each cell for eight curated gene programs using
\texttt{scanpy.tl.score\_genes}. The pseudotime position at which each program's mean
score peaked (across 20 equal-width bins) yielded: AT2 identity at 0.015, AT1
identity at 0.266, apoptosis at 0.591, AT2-to-AT1 differentiation at 0.744,
NF-$\kappa$B at 0.744, oxidative stress at 0.744, interferon at 0.866, senescence
at 0.935 (Supplementary Figure S6; Supplementary Table S8).

The dominant pattern --- homeostatic identity programs peak in the lower third of
pseudotime, injury-associated programs peak in the upper half --- is stable across
all ablations. The fine-grained ordering \textit{within} the injury phase (specifically,
apoptosis peaking before inflammatory/stress programs) did not match our
pre-specified expectation of an interferon-first cascade, and it inverts under
MSigDB-hallmark substitution (Ablation 8, $\rho = -0.80$). We therefore treat the
fine ordering as hypothesis-generating. The identity-decline pattern is a robust
result; the specific injury-phase sequence is not.

\begin{figure}[htbp]
\centering
\includegraphics[width=\textwidth]{results/figures/fig3_program_dynamics.png}
\caption{\textbf{Gene-program dynamics along pseudotime.} Mean program scores across 20 pseudotime bins. Homeostatic identity programs (AT2, AT1) peak early; injury-associated programs peak in the upper half.}
\label{fig:program-dynamics}
\end{figure}

\subsection{Independent replication: geometric findings are portable, annotation-dependent findings are not}

The independent replication cohort (89,736 alveolar AT1/AT2 nuclei from 35
datasets and 618 donors; SCP1219 excluded) provides a structured test of which
findings are portable.

\subsubsection{Portable}

Dispersion replicated strongly (Levene $p \approx 10^{-137}$; variance ratio
1.65). The donor-level direction of the pseudotime-analogue shift replicated
($p = 0.002$; 43 COVID versus 575 control donors). Interferon-response and
AT1-identity programs scored significantly higher in COVID-19 cells at the cell
level ($p \approx 10^{-194}$ and $p \approx 10^{-104}$, respectively).

\subsubsection{Not portable}

Cell-level injury-composite shift failed ($p = 1.0$), mirroring the
primary-cohort centroid-displacement failure. NF-$\kappa$B, oxidative stress,
AT2-to-AT1 differentiation, apoptosis, and senescence program scores did not show
COVID-19 enrichment at the cell level. Marker-threshold DATP enrichment reversed
direction (0.48-fold instead of 3.1-fold).

\subsubsection{Interpretation}

The overall pattern is that \textit{geometric} findings (dispersion, donor-level
direction of shift) are portable across cohorts and methods.
\textit{Annotation-dependent} and \textit{threshold-based} findings (DATP marker
co-expression, specific program peak positions) are atlas-specific. This is a
more precise and more honest view of generalizability than a single-dataset
analysis can provide.

\begin{figure}[htbp]
\centering
\includegraphics[width=\textwidth]{results/figures/fig7_replication_summary.png}
\caption{\textbf{Independent replication summary.} (a) Portability matrix: which findings replicate across 35 datasets and 618 donors. (b) Per-program significance in the replication cohort (cell-level Mann--Whitney). (c) State score distributions by condition in the replication cohort (all $p < 10^{-300}$).}
\label{fig:replication}
\end{figure}

\section{Discussion}

\subsection{Principal finding}

This study set out to distinguish two competing geometric models of alveolar
epithelial injury in lethal COVID-19. Under the \textit{coherent displacement} model,
injury shifts the entire cell population toward a single new location in
transcriptomic space. Under the \textit{dispersion} model, injury fans cells outward
into a heterogeneous cloud. Our data support the dispersion model. Centroid
displacement was not detected ($p = 1.0$), while within-group dispersion was the
single strongest signal in the dataset (variance ratio 2.3; Levene
$p < 10^{-30}$), replicated cross-cohort (variance ratio 1.65;
$p \approx 10^{-137}$), and robust across all tested analytical perturbations. We
describe this pattern as \textit{dispersion-dominated loss of alveolar epithelial state
coherence}: COVID-19 alveolar cells lose their coordinated homeostatic identity
and scatter across a heterogeneous range of injury, repair, and stress states
while maintaining approximately the same average distance from the healthy
centroid.

\subsection{Why centroid displacement failed and why that matters}

The centroid-displacement test was pre-specified because, under a simple intuition
that ``sick cells are further from healthy cells,'' it should have been the most
direct readout. It was not. The median distance of COVID-19 cells from the healthy
centroid (10.68) was actually lower than the median distance of healthy cells from
their own centroid (11.56). This makes biological sense: some COVID-19 cells
remain near homeostatic identity while others occupy distant, heterogeneous
positions. A mean-displacement test averages over these subpopulations and sees
nothing. This negative result is not a failure of the study but a positive
finding: it rules out the simpler ``one coherent damaged state'' model and motivates
the dispersion-based framework.

\subsection{Donor-aware pseudotime shift}

The pseudotime shift (+0.048 on a [0,1] axis; permutation $p = 0.001$) is modest
but robust: preserved in 27/27 leave-one-donor-out iterations, amplified 2.6-fold
under Harmony batch correction, confirmed at reduced magnitude by Palantir, and
replicated at the donor-pseudobulk level in an independent cohort ($p = 0.002$).
The amplification under Harmony is important because it indicates that the
uncorrected estimate is a conservative lower bound --- donor-level batch structure
was \textit{attenuating} the signal, not inflating it.

\subsection{Portable findings versus atlas-specific findings}

The independent replication distinguishes two classes of findings. Geometric
signatures --- within-group dispersion and donor-level direction of the
pseudotime-analogue shift --- are portable across 35 datasets and 618 donors.
Annotation-dependent signatures --- marker-threshold DATP enrichment, specific
program peak positions --- are not. This distinction suggests that the geometric
properties of cell-state distributions (variance, spread, trajectory direction
aggregated to donors) may be more reliable cross-cohort readouts of tissue injury
than the precise identities or thresholds of specific cell populations.

\subsection{Transitional cells: primary-atlas enrichment, not portable thresholds}

The 3.1-fold enrichment of KRT8$^+$/CLDN4$^+$ transitional cells in COVID-19 is
consistent with accumulation of cells that initiated AT2-to-AT1 repair but did
not complete it. However, the marker-threshold definition of this population does
not transfer cross-cohort: the replication cohort shows 47.6\% DATP-positive among
controls versus 22.9\% among COVID-19, the opposite direction. This is not
evidence against the biological reality of transitional cells; it is evidence that
a fixed marker-expression threshold applied across datasets with different
normalization does not produce comparable proportions.

\subsection{Gene-program ordering: stable framework, fragile details}

The stable finding is that homeostatic identity programs (AT2, AT1) decline
monotonically with pseudotime while injury-associated programs rise in the upper
half. The fragile finding is the specific order within the injury phase: it
inverts under MSigDB-hallmark substitution ($\rho = -0.80$) and under
COVID-extreme rooting ($\rho = -0.10$). We do not claim a specific
injury-program cascade.

\subsection{Mechanistic decomposition of the fan-out}

A linear regression of distance-to-healthy-centroid on program scores
($R^{2} = 0.11$) identified senescence (coefficient = 6.83), oxidative stress
(3.83), and NF-$\kappa$B inflammatory (1.95) as the top contributors to the
fan-out pattern. AT2 identity ($-1.13$) and AT1 identity ($-1.02$) carried negative
coefficients, as expected: cells retaining homeostatic identity are closer to the
healthy centroid.

\subsection{Trajectory is alveolar-specific}

Ablation 6 added airway epithelial cells and the pseudotime shift collapsed (MPD
= $-0.008$). The trajectory we describe is a property of the alveolar compartment
in isolation.

\subsection{Limitations}

\begin{itemize}
\item \textit{Cross-sectional design.} All cells were collected at a single time point. Pseudotime is a computational ordering, not a measurement of elapsed time.
\item \textit{Sample size.} The primary atlas contains 27 donors (7 versus 20). Leave-one-donor-out, Harmony, and cross-cohort replication (618 donors) collectively address this, but primary donor-level comparisons are low-powered.
\item \textit{Donor-level modeling.} A donor-level OLS model (median pseudotime $\sim$ condition, $n = 27$) yielded a COVID-19 coefficient of +0.043 (95\% CI [0.010, 0.076]; $p = 0.013$; $R^{2} = 0.22$). The donor-level bootstrap 95\% CI for the median pseudotime difference was [$-0.013$, 0.086]; the CI marginally includes zero.
\item \textit{Lethal cases only.} No mild or moderate COVID-19 cases are available.
\item \textit{Post-mortem effects.} Peri-mortem hypoxia may influence gene expression, particularly stress and apoptosis programs.
\item \textit{Replication scope.} The independent cohort uses an 84-gene panel and cannot support full manifold reconstruction.
\item \textit{Gene-list sensitivity.} The fine-grained injury-program ordering is not robust to gene-list choice.
\item \textit{DATP portability.} Marker-threshold DATP enrichment does not transfer cross-cohort.
\end{itemize}

\subsection{Conclusion}

Lethal COVID-19 does not push alveolar epithelial cells coherently away from the
healthy state. It fans them outward across heterogeneous injury states while
increasing their transcriptomic dispersion, with accumulation of a
KRT8$^+$/CLDN4$^+$ transitional population at intermediate trajectory positions.
The strongest and most portable finding is within-group dispersion, which
replicates across 35 independent datasets and 618 donors. Donor-level trajectory
direction replicates; cell-level displacement does not; and threshold-based
transitional-state definitions are atlas-specific.

\subsection{Future directions}

Five planned extensions would strengthen the generalizability and mechanistic depth of these findings:

\begin{enumerate}

\item \textbf{Cell-level mixed-effects models.} The current donor-level OLS treats donor medians as independent observations. Replacing this with mixed-effects models (pseudotime $\sim$ condition $+$ (1$|$donor\_id) and dispersion $\sim$ condition $+$ (1$|$donor\_id)) would properly account for within-donor correlation and address the pseudoreplication concern inherent in cell-level Levene tests with $p < 10^{-30}$.

\item \textbf{Full-transcriptome independent replication.} The current replication cohort is restricted to an 84-gene panel. Applying the full PCA-to-DPT-to-three-test pipeline to Delorey et~al.\ (2021, GSE171524; $\sim$15 COVID-19 + controls) would provide a complete independent test of the dispersion framework without gene-panel limitations.

\item \textbf{Cross-disease generalization.} Testing whether the dispersion framework extends beyond COVID-19 to idiopathic pulmonary fibrosis using Adams et~al.\ (2020, GSE136831; $\sim$32 IPF + 28 controls) would determine whether loss of state coherence is a general property of alveolar injury or specific to viral pneumonia.

\item \textbf{Cross-species replication with temporal calibration.} Strunz et~al.\ (2020, GSE141259) provides mouse bleomycin lung injury data with a time series (day 0 through day 56) and Sftpc-CreERT2 lineage tracing. This dataset would enable: (a)~cross-species validation of the dispersion framework via ortholog-mapped gene programs, (b)~calibration of computational pseudotime against real elapsed time, and (c)~lineage-trace-based validation of trajectory direction without computational assumptions. The Krt8$^+$ transitional state was first characterized in this dataset, making it the natural ground-truth system for our framework.

\item \textbf{RNA velocity and classifier-based state transfer.} If BAM or loom files are available for SCP1219, RNA velocity would provide direction-independent trajectory validation. Additionally, training a logistic regression classifier on transitional versus non-transitional cells using the 84-gene panel and applying it cross-cohort would convert the DATP marker-threshold failure into a positive result: geometric signatures transfer even when marker thresholds do not.

\end{enumerate}

Together, these extensions would establish a four-layer generalization stack --- human COVID-19, human multi-cohort, human IPF, and mouse bleomycin --- sufficient to position the dispersion framework as a general tool for characterizing epithelial injury geometry across diseases and species.

\section{Methods}

\subsection{Data}

SCP1219 \citep{Melms2021}, 116,313 nuclei from 27 donors. We retained AT1, AT2, and
ECM-high epithelial cells (22,128 cells).

\subsection{Preprocessing}

Normalize to 10,000 counts per cell; $\log(1+x)$; 3,000 HVGs (Seurat-v3); scale
(max 10); 50-component PCA (seed 42).

\subsection{Graph and embeddings}

$k$-NN $k = 15$; UMAP (min\_dist = 0.3); 15-component diffusion map.

\subsection{Batch correction}

Harmony \citep{Korsunsky2019} with \texttt{donor\_id} as batch key, max\_iter\_harmony = 20,
$\sigma = 0.1$. Converged in 3 iterations.

\subsection{Pseudotime}

DPT \citep{Haghverdi2016} via \texttt{scanpy.tl.dpt}, rooted at healthy AT2 centroid in
diffusion space; normalized to [0,1]. Palantir \citep{Setty2019} with 1,200 waypoints
and same root for cross-validation.

\subsection{Gene programs}

Eight programs (AT2 identity; AT1 identity; interferon; NF-$\kappa$B; oxidative
stress; AT2-to-AT1 differentiation; apoptosis; senescence) scored via
\texttt{scanpy.tl.score\_genes} against full normalized expression matrix.

\subsection{Statistical tests}

\begin{itemize}
\item Centroid displacement: one-sided Mann--Whitney U (COVID $>$ Healthy) on per-cell distance to healthy centroid in PCA; effect size rank-biserial $r$.
\item Dispersion: Levene's test on per-cell distance to own-condition centroid.
\item Pseudotime: 1,000-permutation test on median difference.
\item Program ordering: Spearman $\rho$ on peak-position vectors.
\end{itemize}

\subsection{Replication}

89,736 alveolar AT1/AT2 cells from cellxgene census (35 datasets, 618 donors;
SCP1219 excluded). 84-gene panel; per-donor cap 200 cells; injury composite
scored from six non-homeostatic programs. Cell-level and donor-level tests.

\subsection{Software}

Python 3.12.3; Scanpy 1.12.1; AnnData 0.12.10; NumPy 2.4.3; SciPy 1.17.1;
pandas 2.3.3; harmonypy 0.2.0; palantir 1.4.4. Seeds fixed to 42.

\section{Tables}

\begin{table}[htbp]
\centering
\caption{Three pre-specified tests of epithelial state coherence.}
\label{tab:three-tests}
\begin{tabular}{lllll}
\toprule
Test & Statistic & $p$-value & Effect size & Direction \\
\midrule
Mann--Whitney U (displacement) & $5.21 \times 10^7$ & 1.0 & $r = 0.14$ & opposite to alternative \\
Levene's test (dispersion) & 276.5 & $< 10^{-30}$ & var.\ ratio 2.3 & COVID $>$ Healthy \\
Permutation (pseudotime shift) & --- & 0.000999 & MPD +0.048 & COVID $>$ Healthy \\
\bottomrule
\end{tabular}
\end{table}

\begin{table}[htbp]
\centering
\caption{Ten pre-specified ablations summary.}
\label{tab:ablations}
\small
\begin{tabular}{llllp{4cm}}
\toprule
Ablation & MPD & $\rho$ & $r$ & Notes \\
\midrule
1 AT2-only & +0.133 & 0.33 & $-0.29$ & sign flip on $r$ \\
2 UMAP / diffmap & +0.048 & 0.55 & +0.14 & unchanged \\
3 Harmony corrected & +0.125 & 0.45 & +0.14 & amplifies shift \\
3 Uncorrected & +0.048 & 0.55 & +0.14 & baseline \\
4 No apoptosis genes & +0.074 & 0.50 & +0.14 & slight up \\
5 Four roots & $-0.037$ to +0.048 & $-0.10$ to 0.64 & +0.14 & root-dependent \\
6 Broader epithelial & $-0.008$ & 0.83 & +0.25 & \textbf{shift abolished} \\
7 LOO (27 donors) & 0.028--0.202 & varies & varies & direction 27/27 \\
8 Curated-only & +0.048 & 1.00 & +0.14 & stable \\
8 MSigDB-only & +0.048 & $-0.80$ & +0.14 & \textbf{ordering inverted} \\
9 Palantir & +0.028 & 0.48 & +0.14 & alt.\ algorithm \\
10 Permutation & $p=0.001$ & --- & --- & 1000 shuffles \\
\bottomrule
\end{tabular}
\end{table}

\begin{table}[htbp]
\centering
\caption{Composite assessment (primary + replication).}
\label{tab:composite}
\begin{tabular}{lll}
\toprule
Criterion & Primary (SCP1219) & Replication (census) \\
\midrule
Centroid displacement & \textbf{fails} ($p=1.0$) & \textbf{fails cell-level} ($p=1.0$) \\
Dispersion (Levene) & supported ($p \sim 0$) & \textbf{supported} ($p \sim 10^{-137}$) \\
Pseudotime enrichment & supported ($p=0.001$) & supported donor-level ($p=0.002$) \\
Harmony-corrected shift & MPD $\times$2.6 & --- \\
KRT8$^+$/CLDN4$^+$ DATPs & 3.1$\times$ enriched & \textbf{0.48$\times$ --- retracted} \\
\bottomrule
\end{tabular}
\end{table}

\section{Data and Code Availability}

The Columbia/NYP COVID-19 Lung Atlas is available through the Broad Institute
Single Cell Portal under accession SCP1219 \citep{Melms2021}. The cellxgene census
replication cohort was accessed via the cellxgene census API (release 2025-11-08).
Analysis code and configuration are available at the project repository.

\begin{thebibliography}{15}

\bibitem[Adams et~al.(2020)]{Adams2020}
Adams TS, Schupp JC, Poli S, et~al. (2020).
\newblock Single-cell RNA-seq reveals ectopic and aberrant lung-resident cell populations in idiopathic pulmonary fibrosis.
\newblock \textit{Science Advances}, 6(28):eaba1983.

\bibitem[Barkauskas et~al.(2013)]{Barkauskas2013}
Barkauskas CE, Cronce MJ, Rackley CR, et~al. (2013).
\newblock Type 2 alveolar cells are stem cells in adult lung.
\newblock \textit{Journal of Clinical Investigation}, 123(7):3025--3036.

\bibitem[Borczuk et~al.(2020)]{Borczuk2020}
Borczuk AC, Salvatore SP, Seshan SV, et~al. (2020).
\newblock COVID-19 pulmonary pathology: a multi-institutional autopsy cohort from Italy and New York City.
\newblock \textit{Modern Pathology}, 33(11):2156--2168.

\bibitem[Delorey et~al.(2021)]{Delorey2021}
Delorey TM, Ziegler CGK, Heimberg G, et~al. (2021).
\newblock COVID-19 tissue atlases reveal SARS-CoV-2 pathology and cellular targets.
\newblock \textit{Nature}, 595(7865):107--113.

\bibitem[Habermann et~al.(2020)]{Habermann2020}
Habermann AC, Gutierrez AJ, Bui LT, et~al. (2020).
\newblock Single-cell RNA sequencing reveals profibrotic roles of distinct epithelial and mesenchymal lineages in pulmonary fibrosis.
\newblock \textit{Science Advances}, 6(28):eaba1972.

\bibitem[Haghverdi et~al.(2016)]{Haghverdi2016}
Haghverdi L, B{\"u}ttner M, Wolf FA, Buettner F, Theis FJ. (2016).
\newblock Diffusion pseudotime robustly reconstructs lineage branching.
\newblock \textit{Nature Methods}, 13(10):845--848.

\bibitem[Hou et~al.(2020)]{Hou2020}
Hou YJ, Okuda K, Edwards CE, et~al. (2020).
\newblock SARS-CoV-2 reverse genetics reveals a variable infection gradient in the respiratory tract.
\newblock \textit{Cell}, 182(2):429--446.e14.

\bibitem[Kobayashi et~al.(2020)]{Kobayashi2020}
Kobayashi Y, Tata A, Konkimalla A, et~al. (2020).
\newblock Persistence of a regeneration-associated, transitional alveolar epithelial cell state in pulmonary fibrosis.
\newblock \textit{Nature Cell Biology}, 22(8):934--946.

\bibitem[Korsunsky et~al.(2019)]{Korsunsky2019}
Korsunsky I, Millard N, Fan J, et~al. (2019).
\newblock Fast, sensitive and accurate integration of single-cell data with Harmony.
\newblock \textit{Nature Methods}, 16(12):1289--1296.

\bibitem[Melms et~al.(2021)]{Melms2021}
Melms JC, Biermann J, Huang H, et~al. (2021).
\newblock A molecular single-cell lung atlas of lethal COVID-19.
\newblock \textit{Nature}, 595(7865):114--119.

\bibitem[Setty et~al.(2019)]{Setty2019}
Setty M, Kiseliovas V, Levine J, Gayoso A, Mazutis L, Pe'er D. (2019).
\newblock Characterization of cell fate probabilities in single-cell data with Palantir.
\newblock \textit{Nature Biotechnology}, 37(4):451--460.

\bibitem[Strunz et~al.(2020)]{Strunz2020}
Strunz M, Simon LM, Ansari M, et~al. (2020).
\newblock Alveolar regeneration through a Krt8+ transitional stem cell state that persists in human lung fibrosis.
\newblock \textit{Nature Communications}, 11(1):3559.

\bibitem[Wendisch et~al.(2021)]{Wendisch2021}
Wendisch D, Dietrich O, Mari T, et~al. (2021).
\newblock SARS-CoV-2 infection triggers profibrotic macrophage responses and lung fibrosis.
\newblock \textit{Cell}, 184(26):6243--6261.e27.

\bibitem[Wolf et~al.(2018)]{Wolf2018}
Wolf FA, Angerer P, Theis FJ. (2018).
\newblock SCANPY: large-scale single-cell gene expression data analysis.
\newblock \textit{Genome Biology}, 19(1):15.

\bibitem[Ziegler et~al.(2020)]{Ziegler2020}
Ziegler CGK, Allon SJ, Nyquist SK, et~al. (2020).
\newblock SARS-CoV-2 receptor ACE2 is an interferon-stimulated gene in human airway epithelial cells and is detected in specific cell subsets across tissues.
\newblock \textit{Cell}, 181(5):1016--1035.e19.

\end{thebibliography}

\end{document}
